% !TEX program = pdflatex
\documentclass[11pt]{article}
\usepackage[margin=0.8in]{geometry}
\usepackage{amsmath,amssymb}
\usepackage{enumitem}
\usepackage{titlesec}
\usepackage{xcolor}

\titleformat{\section}{\large\bfseries}{}{0em}{}[\vspace{-0.3em}\rule{\textwidth}{0.4pt}]
\titleformat{\subsection}{\normalsize\bfseries}{}{0em}{}
\setlength{\parindent}{0pt}
\setlength{\parskip}{0.4em}

\definecolor{qcolor}{HTML}{1565C0}
\newcommand{\Q}[1]{\textcolor{qcolor}{\textbf{Q: #1}}}

\begin{document}

{\Large\textbf{Week 3 Meeting Cheat Sheet}}\\[0.3em]
Anticipated questions and answers, organized by slide.

\section{General / Big Picture}

\Q{What's the main new result this week?}
The $\alpha\in(2,3)$ gap: no substitution matrix with $|\det M|=1$ can produce $\alpha$ between 2 and 3, but $|\det M|=2$ can. First known 1D example at $\alpha=2.071$.

\Q{How does this fit into the JP?}
This would go in a results section on the substitution matrix parameter space. The gap theorem + $|\det M|>1$ construction is probably the most novel contribution.

\Q{What's the most important open question?}
Whether Silver $\bar\Lambda = 1/4$ exactly. If true, there's a clean rational sequence: lattice$=1/6$, Silver$=1/4$, URL$=1/3$. Proving it would require applying the Zachary--Torquato theta-function method to the Silver chain.

\section{Slide 3: Why spreadability failed for BT}

\Q{Does spreadability ever work for BT?}
Yes, but it needs $t > 10^8$. At $t\in[10^2, 10^5]$ you get $\hat\alpha = 2.26$. At $[10^8, 10^{11}]$ you finally get ${\sim}1.51$. It converges, just very slowly.

\Q{Why does $\bar\Lambda$ from variance work but spreadability doesn't?}
$\bar\Lambda$ is a running average of $\sigma^2(R)$ --- it doesn't need a fitting window or slope extraction. Spreadability requires fitting a power law $\mathcal{E}(t) \sim t^{-(1+\alpha)/2}$, and the correction terms from large $|\lambda_2| = 0.802$ contaminate the fit at moderate $t$.

\Q{Could you fix spreadability by using a larger fitting window?}
You'd need $t > 10^8$, which requires computing the spectral density and spreadability to extremely high precision. Not practical compared to just using variance.

\section{Slide 4: Convergence with $N$}

\Q{How many points do you need for a reliable $\bar\Lambda$?}
About $N \sim 1000$ for substitution chains. The running average stabilizes quickly. This is shown in panel (a) --- both BT and Fibonacci are flat by $N = 10^3$.

\Q{What about for the stealthy patterns?}
Stealthy data came from Sam --- ${\sim}4316$ configurations per $\chi$ value, $N=2000$, $\rho=1$. The ensemble averaging substitutes for the long-chain averaging we do for substitutions.

\section{Slide 5: Exploring cubic substitution matrices}

\Q{Why row sums $\leq 4$?}
Computational tractability. Row sums = number of tiles produced by each substitution rule. Beyond 4, the number of matrices grows fast. With row sums up to 4, there are 34 possible rows, giving $34^3 = 39{,}304$ matrices.

\Q{Why does $|\det M|=1$ matter?}
$\det M$ = product of all eigenvalues. When $|\det M|=1$, the eigenvalues are tightly constrained. If $\lambda_2, \lambda_3$ are a complex conjugate pair: $|\det M| = \lambda_1|\lambda_2|^2 = 1$, which forces $|\lambda_2| = 1/\sqrt{\lambda_1}$, which gives $\alpha = 2$ exactly.

\Q{What's special about complex conjugate eigenvalue pairs?}
For a $3\times3$ matrix, the characteristic polynomial is cubic with real coefficients. If it has one real root ($\lambda_1$) and two complex roots, those must be conjugates. This is the generic case --- having all three real requires special structure.

\Q{Could there be $\alpha\in(2,3)$ with larger row sums?}
Not with $|\det M|=1$. The proof is algebraic and doesn't depend on row sums. For any $|\det M|=1$ matrix: complex pairs give $\alpha=2$ (rank 3) or $\alpha<2$ (rank $\geq4$). All-real cases give $\alpha=3$ in every example we've checked.

\Q{What are the 4,577 remaining matrices?}
They have all-real eigenvalues and all give $\alpha=3$. Like higher-dimensional versions of the metallic-mean chains. No analytical proof for why $\alpha=3$, just numerical confirmation.

\section{Slide 6: Substitution tilings at $\alpha=2$}

\Q{How different are the 504 matrices from each other?}
18 distinct $\lambda_1$ values (range 1.32 to 3.85). Matrices with the same $\lambda_1$ share the same characteristic polynomial but have different substitution rules. $\bar\Lambda$ varies widely (0.26 to 1.7), even among those with the same $\lambda_1$.

\Q{Why only show one example?}
$\bar\Lambda$ was only computed for one representative per $\lambda_1$ group. Computing all 504 is a TODO.

\Q{Is $\alpha=2$ interesting physically?}
It's the boundary between Class I behavior where $\bar\Lambda$ is finite and well-behaved, and the regime where additional eigenvalues degrade hyperuniformity. URL with $a=1$ also has $\alpha=2$.

\section{Slide 7: Metrics for Classes II and III}

\Q{Why not just use $\bar\Lambda$ for all classes?}
$\bar\Lambda = \lim \sigma^2(R)$ as $R\to\infty$. For Class II, $\sigma^2 \sim \Lambda_{II}\ln R$, which grows without bound. For Class III, $\sigma^2 \sim \Lambda_{III} R^{1-\alpha}$. So the limit doesn't exist --- $\bar\Lambda$ is undefined (not infinite, just undefined).

\Q{Why use curve fitting instead of the ratio $\sigma^2/\ln R$?}
The ratio $= \Lambda_{II} + b/\ln R$ converges to $\Lambda_{II}$ from above at finite $R$. This plateau method is biased. The curve fit $\sigma^2 = \Lambda_{II}\ln R + b$ extracts $\Lambda_{II}$ without this bias.

\Q{Are $\Lambda_{II}$ and $\Lambda_{III}$ standard notation?}
No, we defined these. Torquato's papers write out the asymptotic forms but don't name the coefficients. We chose $\Lambda$ with subscripts to be consistent with $\bar\Lambda$ for Class I.

\Q{What's the physical meaning of $\Lambda_{II} = 0.082$?}
It's the rate of logarithmic growth of density fluctuations. Smaller $\Lambda_{II}$ means fluctuations grow more slowly --- closer to Class I behavior.

\section{Slide 8: Updated catalog}

\Q{Why is $\alpha$ ``---'' for lattice and stealthy?}
Their $S(k)$ doesn't follow a power law $S(k)\sim k^\alpha$ near $k=0$. The lattice has Bragg peaks ($S(k)$ is delta functions). Stealthy has $S(k)=0$ by construction for $|k|<K$. Neither fits the power-law framework.

\Q{Can you compare $\bar\Lambda$ across different densities?}
Yes. $\bar\Lambda$ is invariant under rescaling $x_i \to x_i/\rho$. Verified: Fibonacci at $\rho=0.724$ vs.\ rescaled to $\rho=1$, difference $<0.001$.

\Q{Why is stealthy $\bar\Lambda$ so high?}
Stealthy with $\chi=0.1$ has $\bar\Lambda = 1.021$ (not in table). The formula $\bar\Lambda \approx 1/(\pi^2\chi)$ fits well. Small $\chi$ means the $S(k)=0$ constraint is weak, allowing large fluctuations.

\Q{Where does the stealthy data come from?}
Sam's data. About 4316 configurations per $\chi$ value, $N=2000$, $\rho=1$. Disordered ground states found by optimization from random initial positions.

\section{Slides 9--11: Matrices with $2 < \alpha < 3$}

\Q{Is the $|\det M|=1$ result a proof or just numerical evidence?}
For the complex-pair case: it's a proof. $|\det M| = \lambda_1|\lambda_2|^2 = 1$ forces $\alpha=2$. For the all-real case: numerical evidence only (all examples with row sums $\leq4$ give $\alpha=3$).

\Q{Why can't you get $\alpha > 2.2$ with $|\det M|=2$?}
$\alpha = 3 - 2\ln|\det M|/\ln\lambda_1$. With $|\det M|=2$, need $\lambda_1 > 4$. As $\lambda_1$ grows, $\alpha$ approaches 3 but slowly. With $\lambda_1$ around 4--10, $\alpha$ stays in about 2.0--2.2.

\Q{Could you use $|\det M|=3$?}
Need $\sqrt{\lambda_1} > 3$, i.e.\ $\lambda_1 > 9$. That means very large substitution rules. At $\lambda_1=10$: $\alpha = 2.05$. At $\lambda_1=100$: $\alpha = 2.52$. Possible but requires enormous matrices.

\Q{Has anyone else found chains with $\alpha\in(2,3)$?}
Og\u{u}z et al.\ (2019) have a matrix family that implicitly fills this range (their $b=2$ family), but they never noted that $|\det M|=1$ matrices can't reach it, never generated these tilings, and never measured $\bar\Lambda$. The gap observation and the construction are new.

\Q{Is the $\alpha=2.071$ chain hyperuniform?}
Yes. Class I (bounded $\sigma^2$). $|\lambda_2| = 0.449 < 1$, so the condition for hyperuniformity is satisfied.

\section{Slide 12: Silver $\bar\Lambda = 1/4$}

\Q{How confident are you that it's exactly $1/4$?}
$0.2502 \pm 0.0006$, which is $0.3\sigma$ from $1/4$. Statistically consistent, but can't distinguish ``exactly $1/4$'' from ``very close'' numerically. An analytical proof would be needed.

\Q{How would you prove it?}
Zachary \& Torquato (2009) computed Fibonacci $\bar\Lambda = 0.20110$ using a theta-function expansion. The same method applied to Silver could give a higher-precision result or exact value.

\Q{Why would Silver be $1/4$?}
Silver has $\rho = 1/2$ exactly (mean spacing $= 2$), which is unusually clean. If Silver $= 1/4$: lattice $= 1/6$, Silver $= 1/4$, URL $= 1/3$ --- simple rationals.

\Q{Does Fibonacci have a known exact $\bar\Lambda$?}
No. Zachary \& Torquato (2009) report 0.20110, which is higher precision than our 0.201 but still a numerical result. No known closed form.

\section{Slide 13: Metallic means and the URL value $1/3$}

\Q{Why compare to $1/3$?}
As $n$ grows, the metallic-mean chain looks increasingly like a lattice with rare defects. That makes the cloaked URL value at $a=1$, where $\bar\Lambda = (1+1^2)/6 = 1/3$ exactly (Klatt et al.\ 2020), a natural comparison point.

\Q{Is the convergence to $1/3$ proven?}
No. It was a numerical observation from the $n=1$ through 20 data. More recent larger-$n$ runs are noisy and can overshoot $1/3$, so the asymptotic limit should be treated as open rather than established.

\Q{Is there a closed-form formula for $\bar\Lambda(\mu_n)$?}
Open question. Silver might be $1/4$. The relation between the metallic-mean sequence and the URL value $1/3$ is still not understood analytically.

\Q{How is this related to URL?}
URL with parameter $a$ has $\bar\Lambda = (1+a^2)/6$. At $a=0$: perfect lattice ($\bar\Lambda = 1/6$). At $a=1$: maximally disordered ($\bar\Lambda = 1/3$). The metallic-mean data move toward that scale for moderate $n$, but the exact asymptotic relationship is unresolved.

\section{General Technical Questions}

\Q{What's the eigenvalue formula for $\alpha$?}
$\alpha = 1 - 2\ln|\lambda_2|/\ln\lambda_1$, from Og\u{u}z et al.\ (2019). $\lambda_1$ is the largest eigenvalue, $\lambda_2$ is the second largest by absolute value.

\Q{What's the difference between $\alpha$ from $S(k)$ and $\alpha$ from spreadability?}
Same quantity in theory. $S(k) \sim k^\alpha$ as $k\to 0$ defines $\alpha$. Spreadability $\mathcal{E}(t) \sim t^{-(1+\alpha)/2}$ gives the same $\alpha$. In practice, $S(k)$ is hard for substitution tilings (all Bragg peaks, no smooth background). Spreadability works but can be slow. Variance-based $\bar\Lambda$ avoids needing $\alpha$ entirely.

\Q{How do you compute $\sigma^2(R)$?}
Place random windows of half-width $R$ along the chain, count points in each, compute variance. We use 20,000--50,000 windows depending on chain length.

\Q{What system sizes are you using?}
Typically $N = 500$k to 1M for routine computations. High-precision: $N = 3.88$M (Silver confirmation), $N = 9.37$M (Silver figure), $N = 17.8$M ($\alpha=2.071$ chain).

\end{document}
